\section{Conclusion}

We set out to reduce the security workload in continuous security compliance by automatically including security requirements on the backlog items engineers and security experts already use.
The three contributions of this paper address complementary parts of that goal.

C1 provides an independently expert-annotated backlog-security dataset: $288$ items annotated by a pool of nine security practitioners, resulting in $972$ independent binary assessments, Fleiss' $\kappa = 0.787$, and additional Firesmith security-requirement metadata available for many items.
C1 is positioned as a high-precision complement to the larger Wu et al.\ benchmark family~\cite{wu2022data-quality}, and is released publicly to allow downstream work on security-relevance detection to train and audit against expert supervision.

C2 is a recall-oriented binary classifier for security-relevance.
A frozen MPNet encoder paired with an elastic-net logistic-regression head reaches F2 $=0.774$ on the C1 held-out split and, applied zero-shot to the five Wu cleaned benchmarks, achieves cross-distribution G-measure $\approx 0.65$, competitive with the published classical-ML and open-source-GPT baselines despite a prevalence shift from $27.8\%$ on C1 to $1.9$--$17.9\%$ on the Wu projects.

C3 is a four-stage retrieval-augmented generation pipeline that maps a flagged backlog item to the subset of security requirements from relevant security documents.
The C3 retrieval pipeline was evaluated preliminarily on two industrial product backlogs with industrial internal policies as the primary corpus and CIS Benchmarks as a public-standard probe.
Semi-structured practitioner interviews produced a positive preliminary signal on items with well-structured descriptions, identified integration into backlog systems and security requirement summarization as the most important additions, and suggested a consistent lesson about the granularity gap between security requirement documents and backlog items.
Rule-level documents yield more precise, actionable matches than higher-abstraction capability standards.

Taken together, the results provide preliminary evidence that security-relevance can be decided directly from ordinary backlog text to support engineering work through concrete security requirements attached to those items in a way that potentially makes security both more visible and better integrated in the software engineering lifecycle, without requiring engineers to leave their existing workflow.

~\label{sec:future-work}
Open directions follow directly from the preliminary evaluation.
First, a complete evaluation of C3 is required, including more practitioner interviews, multiple backlogs, ideally across organizations, and a controlled study design capable of measuring workflow impact and adoption.
This would be achieved through a longitudinal study.
Second, applying C3 to a second industrial corpus, and ideally to a second organization's policies, would convert the current document-agnostic architectural claim into a measured generalizability result.
Third, in-tracker integration and per-security requirements action-item summarization are concrete engineering steps that the practitioners identified as the dominant preconditions for practical adoption.
